\documentclass[10pt,a4paper]{amsart}
\usepackage{amsmath,amssymb,amsthm,graphicx}
\usepackage{verbatim}
\newenvironment{reference}{\comment}{\endcomment}
%\newenvironment{reference}{}{}
\newenvironment{slogan}{\comment}{\endcomment}
\newenvironment{history}{\comment}{\endcomment}

% For commutative diagrams we use Xy-pic
\usepackage[all]{xy}

% We use 2cell for 2-commutative diagrams.
\xyoption{2cell}
% Theorem environments.
%
\theoremstyle{plain}
\newtheorem{theorem}[subsection]{Theorem}
\newtheorem{proposition}[subsection]{Proposition}
\newtheorem{lemma}[subsection]{Lemma}

\theoremstyle{definition}
\newtheorem{definition}[subsection]{Definition}
\newtheorem{example}[subsection]{Example}
\newtheorem{exercise}[subsection]{Exercise}
\newtheorem{situation}[subsection]{Situation}

\theoremstyle{remark}
\newtheorem{remark}[subsection]{Remark}
\newtheorem{remarks}[subsection]{Remarks}

\numberwithin{equation}{subsection}

% Macros
%
\def\lim{\mathop{\mathrm{lim}}\nolimits}
\def\colim{\mathop{\mathrm{colim}}\nolimits}
\def\Spec{\mathop{\mathrm{Spec}}}
\def\Hom{\mathop{\mathrm{Hom}}\nolimits}
\def\Ext{\mathop{\mathrm{Ext}}\nolimits}
\def\SheafHom{\mathop{\mathcal{H}\!\mathit{om}}\nolimits}
\def\SheafExt{\mathop{\mathcal{E}\!\mathit{xt}}\nolimits}
\def\Sch{\mathit{Sch}}
\def\Mor{\mathop{\mathrm{Mor}}\nolimits}
\def\Ob{\mathop{\mathrm{Ob}}\nolimits}
\def\Sh{\mathop{\mathit{Sh}}\nolimits}
\def\NL{\mathop{N\!L}\nolimits}
\def\CH{\mathop{\mathrm{CH}}\nolimits}
\def\proetale{{pro\text{-}\acute{e}tale}}
\def\etale{{\acute{e}tale}}
\def\QCoh{\mathit{QCoh}}
\def\Ker{\mathop{\mathrm{Ker}}}
\def\Im{\mathop{\mathrm{Im}}}
\def\Coker{\mathop{\mathrm{Coker}}}
\def\Coim{\mathop{\mathrm{Coim}}}

% Boxtimes
%
\DeclareMathSymbol{\boxtimes}{\mathbin}{AMSa}{"02}

%
% Macros for moduli stacks/spaces
%
\def\QCohstack{\mathcal{QC}\!\mathit{oh}}
\def\Cohstack{\mathcal{C}\!\mathit{oh}}
\def\Spacesstack{\mathcal{S}\!\mathit{paces}}
\def\Quotfunctor{\mathrm{Quot}}
\def\Hilbfunctor{\mathrm{Hilb}}
\def\Curvesstack{\mathcal{C}\!\mathit{urves}}
\def\Polarizedstack{\mathcal{P}\!\mathit{olarized}}
\def\Complexesstack{\mathcal{C}\!\mathit{omplexes}}
% \Pic is the operator that assigns to X its picard group, usage \Pic(X)
% \Picardstack_{X/B} denotes the Picard stack of X over B
% \Picardfunctor_{X/B} denotes the Picard functor of X over B
\def\Pic{\mathop{\mathrm{Pic}}\nolimits}
\def\Picardstack{\mathcal{P}\!\mathit{ic}}
\def\Picardfunctor{\mathrm{Pic}}
\def\Deformationcategory{\mathcal{D}\!\mathit{ef}}
\usepackage{cleveref}
\crefname{originallemma}{original lemma}{original lemmas}
\Crefname{originallemma}{Original lemma}{Original lemmas}
\makeatletter
\newlabel{item-f-ideal}{{00HG}{source}{}{sourcepointer-1}{}}
\newlabel{item-f-ideal@cref}{{[originallemma][0][]00HG}{[1][0][]source}{}{}{}}
\makeatother
\crefname{originallemma}{original lemma}{original lemmas}
\Crefname{originallemma}{Original lemma}{Original lemmas}
\makeatletter
\newlabel{lemma-flat}{{00HD}{source}{}{sourcepointer-2}{}}
\newlabel{lemma-flat@cref}{{[originallemma][0][]00HD}{[1][0][]source}{}{}{}}
\makeatother
\crefname{originallemma}{original lemma}{original lemmas}
\Crefname{originallemma}{Original lemma}{Original lemmas}
\makeatletter
\newlabel{lemma-flat-base-change}{{00HI}{source}{}{sourcepointer-3}{}}
\newlabel{lemma-flat-base-change@cref}{{[originallemma][0][]00HI}{[1][0][]source}{}{}{}}
\makeatother
\crefname{originallemma}{original lemma}{original lemmas}
\Crefname{originallemma}{Original lemma}{Original lemmas}
\makeatletter
\newlabel{lemma-flat-localization}{{00HT}{source}{}{sourcepointer-4}{}}
\newlabel{lemma-flat-localization@cref}{{[originallemma][0][]00HT}{[1][0][]source}{}{}{}}
\makeatother
\crefname{originallemma}{original lemma}{original lemmas}
\Crefname{originallemma}{Original lemma}{Original lemmas}
\makeatletter
\newlabel{lemma-local-flat-ff}{{00HR}{source}{}{sourcepointer-5}{}}
\newlabel{lemma-local-flat-ff@cref}{{[originallemma][0][]00HR}{[1][0][]source}{}{}{}}
\makeatother
\title{Flatness of a module and of its base change along a local diagram, both directions}
\author{The Stacks Project Authors}
\date{stacks-project:algebra.tex:local-snapshot-20261012; GNU FDL 1.2 or later}
\begin{document}
\maketitle
\noindent\textit{Editorial note.} All delivered mathematical text is the author's own, copied from the pinned Stacks Project chapter file algebra.tex; no statement or proof step was rewritten, reordered or removed.\par\medskip
\noindent\textit{Editorial note.} Preamble provenance, as required for this source: the frozen snapshot stages the chapter's own preamble file beside it (preamble.tex, 217 lines, 8116 bytes, the same shared Stacks-wide preamble the published chow and etale-cohomology cards use); the first line of the pinned chapter file is \textbackslash{}input\{preamble\}, so no preamble text lives inside the chapter itself. The card takes that file's lines 8-18 (the verbatim package and the three \textbackslash{}comment-based environments reference/slogan/history at lines 9-12, section 13.2, together with the Xy-pic setup) and its lines 154 to the runtime end of the file (the native theorem declarations and mathematical macros). Not one character was changed, no package was added, no macro was defined, and the excerpt serves only as preamble, never as a span coordinate.\par\medskip
\noindent\textit{Editorial note.} The same-chapter pointers are located at their \textbackslash{}label lines inside the pinned chapter file; the cross-chapter pointers are located in the project's own tag list staged in the snapshot (sources/stacks-chow/tags.txt), at the verbatim lines carrying their canonical tags.\par\medskip
\noindent\textit{Editorial note.} The counted statement contains an Xy-pic square and the proof refers to several same-chapter lemmas by label; the references are recorded as navigation stubs and the Xy-pic package is provided by the preamble spans.\par\medskip
\noindent\textit{Editorial note.} The counted spans contain no \textbackslash{}cite, so the card carries no bibliography.\par\medskip
\noindent\textit{Allowed standard prerequisites (editorial).}\begin{itemize}
\item Preservation of flatness by base change. Invoked in the counted proof at lines 24426-24428. Gives the first direction directly and is reused after localising.
\item Flatness is preserved by localisation. Used in the counted proof at lines 24429-24431. Identifies the base-changed module with a localisation of the tensor product, which is the form in which flatness is known.
\end{itemize}\medskip
\ifdefined\hypertarget\hypertarget{sourcepointer-1}{}\fi
\noindent\textit{Source reference.} Original lemma 00HG: sources/stacks-algebra/pinned-original-source.tex, source lines 8936--8936. Navigation pointer for \textbackslash{}ref\{item-f-ideal\} used inside the counted span (line 24441); the argument that consumes it is retained in the span.\par\medskip
\ifdefined\hypertarget\hypertarget{sourcepointer-2}{}\fi
\noindent\textit{Source reference.} Original lemma 00HD: sources/stacks-algebra/pinned-original-source.tex, source lines 8925--8925. Navigation pointer for \textbackslash{}ref\{lemma-flat\} used inside the counted span (line 24441); the argument that consumes it is retained in the span.\par\medskip
\ifdefined\hypertarget\hypertarget{sourcepointer-3}{}\fi
\noindent\textit{Source reference.} Original lemma 00HI: sources/stacks-algebra/pinned-original-source.tex, source lines 9106--9106. Navigation pointer for \textbackslash{}ref\{lemma-flat-base-change\} used inside the counted span (line 24428); the argument that consumes it is retained in the span.\par\medskip
\ifdefined\hypertarget\hypertarget{sourcepointer-4}{}\fi
\noindent\textit{Source reference.} Original lemma 00HT: sources/stacks-algebra/pinned-original-source.tex, source lines 9413--9413. Navigation pointer for \textbackslash{}ref\{lemma-flat-localization\} used inside the counted span (line 24432); the argument that consumes it is retained in the span.\par\medskip
\ifdefined\hypertarget\hypertarget{sourcepointer-5}{}\fi
\noindent\textit{Source reference.} Original lemma 00HR: sources/stacks-algebra/pinned-original-source.tex, source lines 9401--9401. Navigation pointer for \textbackslash{}ref\{lemma-local-flat-ff\} used inside the counted span (line 24439); the argument that consumes it is retained in the span.\par\medskip
\begin{lemma}
\label{lemma-base-change-flat-up-down}
Let
$$
\xymatrix{
S \ar[r] & S' \\
R \ar[r] \ar[u] & R' \ar[u]
}
$$
be a commutative diagram of local homomorphisms of local rings.
Assume that $S'$ is a localization of the tensor product $S \otimes_R R'$.
Let $M$ be an $S$-module and set $M' = S' \otimes_S M$.
\begin{enumerate}
\item If $M$ is flat over $R$ then $M'$ is flat over $R'$.
\item If $M'$ is flat over $R'$ and $R \to R'$ is flat then
$M$ is flat over $R$.
\end{enumerate}
In particular we have
\begin{enumerate}
\item[(3)] If $S$ is flat over $R$ then $S'$ is flat over $R'$.
\item[(4)] If $R' \to S'$ and $R \to R'$ are flat then $S$ is flat over $R$.
\end{enumerate}
\end{lemma}

\begin{proof}
Proof of (1). If $M$ is flat over $R$, then $M \otimes_R R'$
is flat over $R'$ by
Lemma \ref{lemma-flat-base-change}.
If $W \subset S \otimes_R R'$ is the multiplicative subset such that
$W^{-1}(S \otimes_R R') = S'$ then $M' = W^{-1}(M \otimes_R R')$.
Hence $M'$ is flat over $R'$ as the localization of a flat module, see
Lemma \ref{lemma-flat-localization} part (5). This proves (1) and in
particular, we see that (3) holds.

\medskip\noindent
Proof of (2). Suppose that $M'$ is flat over $R'$ and $R \to R'$ is flat.
By (3) applied to the diagram reflected in the northwest diagonal
we see that $S \to S'$ is flat. Thus $S \to S'$ is faithfully flat by
Lemma \ref{lemma-local-flat-ff}.
We are going to use the criterion of
Lemma \ref{lemma-flat} (\ref{item-f-ideal})
to show that $M$ is flat.
Let $I \subset R$ be an ideal. If $I \otimes_R M \to M$ has a kernel,
so does $(I \otimes_R M) \otimes_S S' \to M \otimes_S S' = M'$.
Note that $I \otimes_R R' = IR'$ as $R \to R'$ is flat, and that
$$
(I \otimes_R M) \otimes_S S' =
(I \otimes_R R') \otimes_{R'} (M \otimes_S S') =
IR' \otimes_{R'} M'.
$$
From flatness of $M'$ over $R'$
we conclude that this maps injectively into $M'$.
This concludes the proof of (2), and hence (4) is true as well.
\end{proof}

\end{document}
